\documentclass{article}
\usepackage{amsmath,amssymb}
\begin{document}
\section*{Native unit and binary action: audit and boundary}
Issue 105; base 223, commit \texttt{6933abe13d0cc79fff9ef28a23eacc32b7afc68d}.
The TeX source \emph{One relative construction for context, content and effects},
blob \texttt{b1e84592f59558c851faabb5291d7392d2b2f908}, section Internal
constructors, labels \texttt{eq:groth}, \texttt{eq:grothmonad},
\texttt{prop:groth}, works in a 2-category and generally yields a pseudomonad.
Here $B$ is an ordinary category, the indexed grades and all presheaves
are strict, and the supplied endpoint presentation $P$ is split. The
following statements concern actual functors and natural transformations.
\paragraph{Global unit.} Native inputs: $I_A$, $\epsilon_A:e_AI_A\to A$,
$a^*I_A=I_D$, and equality of endpoint pastes over $a:D\to A$, including
the grade transport. None mentions $Y$. \texttt{ContextadUnit.unit} gives
$1\Rightarrow L$ on the category of strict Cat-valued presheaves for the
actual Grothendieck step $L$. On objects this is $(I_A,Y(\epsilon_A)y)$.
The grade arrow is identity; the payload is the image of the given arrow
under $Y(\epsilon_A)$ up to its prescribed dependent target. Naturality in
both exterior base and payload transformation is equality on whole functors.
Neither invertibility of discharge nor a target unit is assumed. This is
not the free telescope length-one generator $L\Rightarrow T$.
\paragraph{Binary.} At each $A$, the input is a tensor functor on the
actual source fibre of native two-step grade chains and a natural
$\delta:\otimes;e_A\Rightarrow e_{\rm last}$. The grade category keeps
both changing grades and their induced moving intermediate boundary. The
output \texttt{ContextadTensor.binary} has actual domain $L^2Y(A)$, not a
quotient of its objects. A strict isomorphism of realization charts and
\texttt{read} identify its action on all independent decorated squares.
On $(P,(R,y))$ the output is $(P\otimes R,Y(\delta_{P,R})y)$.
\texttt{rowsGrade} tensors the native grade arrow; \texttt{rowsPayload}
applies $Y(\delta_{\rm source})$ to its transported terminal payload arrow
and then the dependent correction from $\delta$-naturality.
\texttt{binary\_natural} is naturality in $Y$ on full fibre functors.
Exterior naturality across bases is \emph{not} inferred:
\texttt{baseChange\_iff} reflects precisely its missing square to the
native decorated rows, not merely their objects.
\paragraph{Positive exterior result.} \texttt{ContextadTensorLift.Data}
imposes two presheaf-independent component laws on chosen lifts: equality
of reindexed accumulated grades and equality of the two $\delta$/endpoint
pastes (including grade transport). These imply exterior base change on
every actual nested object and on the entire payload-coordinate functor
at each fixed two-step chain. Naturality on changing-grade arrows is
neither assumed nor concluded there.
\paragraph{Representables.} \texttt{ContextadRepresentable} specializes
to discrete Yoneda: $f$ is sent to $\epsilon_A;f$ or $\delta;f$, respectively,
and a payload over a specified grade arrow $\alpha$ exists exactly when
$f=e_A(\alpha);g$. These are component computations for \texttt{prop:groth},
not its local equivalence nor an equivalence between the whole presheaf
Kleisli bicategory and its representable part.
\paragraph{Witnesses and sharpness.} The first probe uses a presheaf of
coordinate-indexed word arrows. The noninvertible unit compares $2\to5$;
word order is distinguished at retained positions, but difference at
coordinate 4 disappears. In a varying-cap tensor $\min(r,A)$ with
noninvertible $\delta$, a square moves the outer and inner grades
$1\to4$, $5\to7$; parallel payloads remain distinguished in the actual
nested action. Another square yields noninvertible output grade $1\to2$.
Exterior restriction $2\to4$ disagrees even on the output grade of a
row ending in 5; a fixed cap, in contrast, satisfies the lift laws.
The second probe is sharper: $K_A=\mathbb N\times W$ for the one-object
free word category $W$, and $e_A$ forgets the word component. The tensor
returns $(r,*)$, with $\delta$ identity, and maps $i\le j$ to $j-i$ copies
of letter $A$. Differences add under composition. Both native lift laws
hold; nevertheless the same moving grade arrow $1\to2$ gives words
$[2]$ versus $[4]$ on restriction along $2\to4$. This violates strict
base change without violating even one fixed-coordinate or object law.
Neither fixture is claimed to be a contextad satisfying the missing law.
\paragraph{Residual and consumer.} Construct the source's cartesian
compatibility of tensor on arbitrary grade arrows and use the native
$\delta$ square to prove the reflected row square. Then assemble global
$L^2\Rightarrow L$ and derive unit/associativity from grade and endpoint
coherence. Only then may 223's universal morphism extend $1_L:L\Rightarrow
L$ to $T\to L$ for an independently constructed native strict target.
No global multiplication, general pseudomonad, equivalence of strict
Tel(span)-algebras with contextads, or issue closure is claimed.
\end{document}
